% ============================================
% Proof of Theorem 1: WAIT Algorithm Optimality
% ============================================

\section{Proof of Theorem~\ref{thm:wait_heavy_traffic}}\label{appexidx::proof of wait_heavy_traffic}

This proof establishes asymptotic optimality of the WAIT algorithm when output lengths are known. The argument proceeds in two parts:
\begin{enumerate}
    \item \textbf{Single-type case}: We construct a coupled dominating process using the Lindley recursion and apply Kingman's bound for queue length expectations.
    \item \textbf{Multiple-type case}: We construct an auxiliary embedded full-threshold process and couple it to event-driven WAIT type by type. This embedded process separates stochastic entry accumulation from the event-driven batch composition while preserving the service opportunities created by the WAIT thresholds.
\end{enumerate}
Key techniques include: Lindley recursion for queue length bounds \citep[Proposition 6.3]{asmussen2003applied}, coupling construction to dominate the original process, Kingman's bound for negative drift processes \citep{kingman1962queues}, and standard random walk results \citep{strait1974maximum}.

Throughout this appendix, we use \(b\) to denote the index of the embedded full-threshold review epochs, which is distinct from the stage index \(s \in \{0,1,\ldots,l_j'\}\) used in the main text for decode stages. We use \(t\) for continuous time and \(B\) for the number of embedded review intervals.

We first consider the single-type case, then extend to multiple types.

\subsubsection{Single-Type Case}

Consider a single prompt type with threshold $n$. For general scheduling algorithms, analyzing this system in continuous time is challenging: KV cache grows dynamically during decode, arrivals are stochastic, and eviction risks complicate the dynamics.

The WAIT policy's \emph{threshold mechanism achieves a dimensionality reduction}. By waiting until exactly $n$ prompts accumulate before initiating a batch, the algorithm enforces a fixed batch size, which ensures constant processing time $\Delta T$ per batch. This decouples the continuous-time memory-constrained dynamics into a tractable discrete-time process: a memory-constrained queueing system becomes a simple random walk. The reduction is enabled by the algorithm design and is what makes the subsequent drift analysis apply.

Specifically, let \(\tilde{\lambda}\) denote the continuous-time arrival rate and let \(\mu=\tilde{\lambda}\Delta T\) be the expected number of arrivals during a threshold-sized batch. The boundary case is \(\mu=n\). In this case the embedded process observed at batch-completion epochs, indexed by \(b\in\mathbb{N}\), is a zero-drift threshold queue. The throughput loss arises from ``stuck'' iterations where insufficient prompts are available.

\begin{lemma}[Queue Length and Stuck Time]\label{lemma::throughput gap, single}
This lemma connects queue accumulation to throughput loss through stuck iterations in the critical case \(\mu=n\). Let \(X^b\) denote the number of arrivals during batch \(b\), where \(X^b \sim \mathrm{Poisson}(n)\). Let \(W^b\) denote the queue length after batch \(b\), with \(W^0=0\). The state transition is:
$$W^{b+1}=W^b+X^b-n\cdot \mathbf{1}\{W^b+X^b\geq n\}.$$
Define \(B_{\emph{stuck}}\) as the number of stuck iterations, where \(W^b+X^b<n\). Then:
$$\mathbb{E}[W^B]= n\cdot\mathbb{E}[B_{\emph{stuck}}].$$
\end{lemma}
\noindent The proof is in Appendix~\ref{appendix_lemma::proof of throughput gap, single}. This lemma connects queue accumulation to throughput loss: stuck iterations directly translate to waiting prompts. To bound $\mathbb{E}[W^B]$, we construct a coupled dominating process that admits tractable analysis via classical queueing techniques.

\begin{lemma}[Coupled Dominating Process]\label{lem:coupled_process, single}
The coupled process starts with a safety buffer (\(2n\) versus 0) and maintains dominance throughout. Define \(\{\tilde{W}^b\}\) by:
$$
\tilde{W}^0 = 2n, \quad \tilde{W}^{b+1} = \max\{ 2n, \tilde{W}^b + X^b - n \}.
$$
Then \(\tilde{W}^b \geq W^b + n\) for all \(b\in \mathbb{N}\).
\end{lemma}
\noindent The proof is in Appendix~\ref{appendix_lemma::proof of coupled_process, single}.

\paragraph{Lindley recursion.} To analyze the coupled process, we apply the Lindley recursion representation \citep[Proposition 6.3]{asmussen2003applied}:
\begin{lemma}[Lindley Representation]\label{lem:lindley_rep}
Let \(S^k=\sum_{r=1}^k (X^r-n)\) be the partial sum process. Then:
$$\tilde{W}^B=2n+\max(S^B, S^B-S^1,\ldots,S^B-S^{B-1},0).$$
\end{lemma}
\noindent This classical queueing representation is made applicable by the WAIT policy's threshold design. Define \(\xi^k = X^k-n\). We use time-reversal symmetry to simplify the maximum expression:
$$\max(S^B, S^B-S^1,\ldots,S^B-S^{B-1},0)=\max_{1\leq k\leq B}\Big(\sum_{r=k}^{B}\xi^r\Big)^+\overset{d}{=}\max_{1\leq k\leq B}\Big(\sum_{r=1}^{k}\eta^r\Big)^+, \quad \eta^r\overset{d}{=}X^1-n.$$
Thus \(\tilde{W}^B \overset{d}{=}2n+\max_{1\leq k\leq B}(\sum_{r=1}^{k}\eta^r)^+\).

Since \(\tilde{W}^B \geq W^B + n\), we have \(\mathbb{E}[W^B]\leq \mathbb{E}[\tilde{W}^B]-n\). Let \(Y^B = \max_{1\leq k\leq B}(\sum_{r=1}^{k}\eta^r)^+\). By \cite[Lemma 2]{strait1974maximum}:
\begin{lemma}[Random Walk Maximum]\label{lem:Y_T}
$$\ex{}{Y^B}=\ex{}{\max_{1\leq k\leq B}(S^k)^+}=\sum_{k=1}^B \frac{1}{k}\ex{}{(S^k)^+}.$$
\end{lemma}
\noindent This standard result \citep{strait1974maximum} decomposes the maximum into a sum over batch indices. For \(\mathbb{E}[(S^k)^+]\), the Cauchy-Schwarz inequality and \(\mathrm{Var}(S^k)=nk\) imply \(\mathbb{E}[(S^k)^+] \leq \sqrt{nk}\). Hence:
$$\mathbb{E}[Y^B] \leq \sqrt{n} \sum_{k=1}^B k^{-1/2}=O(\sqrt{nB}).$$

\begin{lemma}[Throughput Gap Bound]\label{lemma::throughput gap}
$$\mathbb{E}[W^B]-n\leq \mathbb{E}[Y^B]\leq  \sqrt{n} \sum_{k=1}^B \frac{1}{\sqrt{k}}=O(\sqrt{nB}).$$
\end{lemma}
\noindent By Lemma~\ref{lemma::throughput gap, single}, the expected number of stuck iterations is \(\mathbb{E}[B_{\emph{stuck}}]=\mathbb{E}[W^B]/n=O(\sqrt{B})\), so the normalized throughput loss from stuck iterations is \(O(B^{-1/2})\). The same bound applies to every prefix \(b\le B\), and therefore
\[
    \frac{1}{B}\sum_{b=1}^B \mathbb{E}[W^b] \leq \frac{1}{B}\sum_{b=1}^B O(\sqrt{b}) = O(\sqrt{B}).
\]
This time-averaged backlog bound is the quantity used for the latency and TTFT estimates.

\paragraph{Scaling with the horizon.} Let \(\Delta^\pi\) denote the fixed processing time of a threshold-sized batch under policy \(\pi\). In the scaled system, the corresponding processing time is \(\Delta^\pi/\zeta\). Over a calendar horizon \(T\), set
\[
    B_\zeta=\left\lfloor \frac{\zeta T}{\Delta^\pi}\right\rfloor .
\]
Then \(B_\zeta=\zeta T/\Delta^\pi+O(1)\). The terminal backlog bound gives normalized throughput loss \(O(B_\zeta^{-1/2})=O((\zeta T)^{-1/2})\). For delay, the theorem reports service-normalized time, so one scaled threshold-batch interval has length \(\Delta^\pi\); the time-averaged backlog bound therefore gives latency and TTFT contributions \(O(B_\zeta^{1/2})=O((\zeta T)^{1/2})\). The bounded number of service stages is lower order. This yields the single-type throughput, latency, and TTFT orders stated in the theorem.

\subsubsection{Multiple-Type Case}

We now extend to \(m\) known output-length classes. The additional issue is that event-driven WAIT need not process the full class composition in every realized batch: if only a subset of class queues has reached its threshold, only that subset is batched. The proof handles this by introducing an auxiliary embedded full-threshold process observed at deterministic review epochs. The actual policy is the event-driven WAIT policy from the main text, which batches all currently eligible types whenever the server becomes idle. The embedded process is not used to compare the total number of realized batches or the total fixed overhead paid by the two systems. Its role is to put each type-\(j\) threshold queue on a common full-threshold service scale, obtain a Lindley-type recursion, and then transfer the resulting bounds back to the actual policy through a sample path coupling of service opportunities.

For a threshold vector \(\pi=[n_1,\ldots,n_m]\), let
\begin{equation}\label{eq:batch_time_def}
\Delta^\pi \equiv \Delta T_{[1,\ldots,m]}
    =  d_0 + d_1 M^\pi,\qquad
M^\pi = \sum_{j=1}^m n_j(l_j^\prime+1)\left( l_j+\frac{l_j^\prime}{2} \right).
\end{equation}
This is the processing time of the full threshold composition, i.e., the batch that contains \(n_j\) prompts from every decode stage of every type \(j\). The load-balance and memory conditions are
\begin{equation}\label{eq:policy_constraints}
\begin{aligned}
&\lambda_j\Delta^\pi \leq n_j \quad \text{for all } j \in [m], \\
&M^{*}\leq M^{\pi}\leq C.
\end{aligned}
\end{equation}
The first line says that, over one full-threshold iteration, the expected type-\(j\) arrivals do not exceed the \(n_j\) type-\(j\) prompts that can be advanced from each active stage. The second line is the corresponding feasibility condition: the threshold composition is at least on the scale of the fluid memory requirement but does not exceed the physical memory \(C\).

\begin{lemma}[GPU-resident memory invariant]\label{lem:wait_resident_memory_invariant}
Starting from the empty system, WAIT maintains
\[
    N_{js}(t)\le n_j,\qquad j\in[m],\; s=1,\ldots,l_j^\prime,
\]
where \(N_{js}(t)\) is the number of resident type-\(j\) prompts at decode stage \(s\). Stage-0 prompts that have not been selected for prefill have no resident KV cache.
\end{lemma}
\noindent The proof is by induction over batch completions. If type \(j\) is not served in a realized batch, its GPU-resident prompts at decode stages are unchanged. If type \(j\) is served, the invariant before the batch implies that each decode stage contains at most \(n_j\) prompts, so WAIT selects all prompts currently present in that stage. After the batch completes, stage \(s+1\) contains only the prompts advanced from stage \(s\), again at most \(n_j\). Thus no decode stage can exceed \(n_j\) resident prompts.

The invariant gives the memory-feasibility statement needed by the policy. If \(a_{j0}(t)\le n_j\) denotes the number of type-\(j\) prompts selected for prefill in the current batch, then at any decision epoch,
\[
    \sum_{j=1}^m\sum_{s=1}^{l_j^\prime}N_{js}(t)(l_j+s)
    +
    \sum_{j=1}^m a_{j0}(t)l_j
    \le
    \sum_{j=1}^m n_j\sum_{s=0}^{l_j^\prime}(l_j+s)
    =
    M^\pi
    \le C .
\]
Thus unselected resident prompts are included in the memory accounting, while unselected stage-0 prompts remain outside the GPU KV-cache state.

For any realized event-driven WAIT batch \(r\), let \(a^{(r)}_{js}\le n_j\) be the number of selected type-\(j\) prompts at decode stage \(s\), and let \(J_r\) be the set of selected types. The selected batch composition has size
\[
    M_r
    =
    \sum_{j\in J_r}\sum_{s=0}^{l_j^\prime} a^{(r)}_{js}(l_j+s)
    \le
    \sum_{j=1}^m n_j\sum_{s=0}^{l_j^\prime}(l_j+s)
    =
    M^\pi.
\]
Therefore the physical processing time of the realized batch in the \(\zeta\)-scaled system is
\[
    S_r^{(\zeta)}
    =
    \frac{d_0+d_1M_r}{\zeta}
    \le
    \frac{\Delta^\pi}{\zeta}.
\]
This bound concerns processing time, while Lemma~\ref{lem:wait_resident_memory_invariant} gives memory feasibility. Actual batches may contain fewer types and have shorter duration, but no realized WAIT batch is larger than the full-threshold composition. The fixed overhead \(d_0\) is therefore handled through a calendar-time coupling of service opportunities, not by matching realized batches one for one. If a type is absent from a realized batch, then either its entry queue has not yet reached threshold, or the threshold is reached while the server is processing another batch. In the latter case the waiting time until the next review is bounded by the remaining processing time of that in-service batch, which is at most one full-threshold interval.

In the \(\zeta\)-scaled system, one full-threshold review interval has physical length \(\Delta^\pi/\zeta\). Define deterministic review epochs
\[
    u_b=b\Delta^\pi/\zeta,\qquad b=0,1,2,\ldots .
\]
The embedded process reviews the entry queues only at these epochs. We use an end-of-interval accounting convention: arrivals accumulate over \((u_b,u_{b+1}]\), and at \(u_{b+1}\) the embedded process removes one type-\(j\) entry batch of size \(n_j\) whenever the type-\(j\) entry queue has reached \(n_j\); it also advances up to \(n_j\) resident type-\(j\) prompts from each active decode stage. Equivalently, one can view the selected work as occupying the following full review interval. A batch removed at review \(u_{b+1}\) can therefore shift completion accounting by at most \(l_j^\prime+1\) embedded intervals for type \(j\), which contributes only \(O((\zeta T)^{-1})\) to normalized throughput and \(O(1)\) to service-normalized average delay. The full-threshold interval length is feasible by the resident memory invariant above, and it is deliberately chosen as the service scale at which the type-wise threshold queues have Lindley recursions.

Throughout this argument, \(b\) indexes embedded full-threshold review intervals, not realized event-driven WAIT batches.

\begin{lemma}[Type-wise sample path coupling of service opportunities]\label{lem:wait_one_slot_dominance}
Couple actual event-driven WAIT and the auxiliary embedded full-threshold process on the same arrival sample path, both starting empty. Let \(L_\zeta=\Delta^\pi/\zeta\) and \(u_b=bL_\zeta\). For type \(j\), let \(R_j^\pi(t)\) be the cumulative number of type-\(j\) threshold batches of size \(n_j\) started by actual WAIT by time \(t\), and let \(R_j^{\bar\pi}(b)\) be the cumulative number of type-\(j\) batches removed from the embedded entry queue at reviews \(u_1,\ldots,u_b\). Then
\[
    R_j^\pi(u_{b+1})\ge R_j^{\bar\pi}(b),
    \qquad b\ge0,\; j\in[m].
\]
The same one-interval delayed dominance holds for type-\(j\) stage advancements.
\end{lemma}
\noindent To see this, fix a type \(j\) and index its threshold batches in order of eligibility. Let \(\tau_{jq}\) be the first time at which the \(q\)-th type-\(j\) entry batch of size \(n_j\) is present in the actual entry queue, and let \(\sigma_{jq}\) be the time at which actual WAIT starts that batch. Since actual WAIT batches all eligible types whenever the server is idle, the batch can wait after \(\tau_{jq}\) only if another batch is already in service. Every such in-service batch has duration at most \(L_\zeta\) by the bound above. Hence, on every sample path,
\[
    \sigma_{jq}\le \tau_{jq}+L_\zeta .
\]
Now suppose the embedded process has removed \(q\) type-\(j\) batches by review epoch \(u_b\). Because the two systems see the same type-\(j\) arrival process, the \(q\)-th actual type-\(j\) threshold batch has become eligible by time \(u_b\), so \(\tau_{jq}\le u_b\). The preceding inequality gives \(\sigma_{jq}\le u_b+L_\zeta=u_{b+1}\). Therefore actual WAIT has started at least as many type-\(j\) threshold batches by \(u_{b+1}\) as the embedded process has removed by \(u_b\), which proves the stated dominance. This argument is in calendar time and includes the fixed overhead paid by any subset batch.

The stage-advancement statement follows from the same sample path coupling of service opportunities. Use the number of type-\(j\) service opportunities as the induction variable. After \(r\) such opportunities, any type-\(j\) batch that has entered service has advanced at least \(\min\{r,l_j'\}\) decode stages, up to the same one-interval lag. Since a type-\(j\) service opportunity advances a threshold batch from every active type-\(j\) stage, the entry dominance propagates stage by stage and gives the stated dominance for completions.

Lemma~\ref{lem:wait_one_slot_dominance} is the structural step that transfers the embedded-process queue bounds back to event-driven WAIT. It does not require the actual system to use fewer batches than the embedded process. Actual WAIT may realize different subset compositions across time, but for each type the threshold batch removed at an embedded review epoch has a corresponding actual service opportunity within one full-threshold review interval. Thus it is enough to bound the embedded comparison queues: the one-interval lag and the bounded in-service pipeline contribute only \(O((\zeta T)^{-1})\) to normalized throughput and \(O(1)\) to the service-normalized delay metrics.

For the embedded process, let \(\bar W^b_{(j)}\) be the residual type-\(j\) entry queue after the \(b\)-th review, and let
\[
    \bar X^b_{(j)}
    =A_j(u_b,u_{b+1}]
    \sim \text{Poisson}\!\left(\zeta\lambda_j\frac{\Delta^\pi}{\zeta}\right)
    =\text{Poisson}(\lambda_j\Delta^\pi).
\]
The arrivals \(\{\bar X^b_{(j)}\}\) are independent across embedded intervals and do not depend on other type queues. With the end-of-interval accounting convention above, the residual entry queue evolves as
\[
    \bar W^{b+1}_{(j)}
    =
    \bar W^b_{(j)}+\bar X^b_{(j)}
    -
    n_j\,\mathbf{1}\{\bar W^b_{(j)}+\bar X^b_{(j)}\ge n_j\}.
\]
This is the point at which the batch-composition issue is reduced to a stochastic comparison: the only coupling across types is the deterministic feasibility condition \(M^\pi\le C\), while the embedded process gives each type its own threshold queue.

\begin{lemma}[Embedded comparison queue bound]\label{lemma::multiple-type coupled dominace}
For each type \(j\), define
\[
\widetilde W_{(j)}^0=2n_j,\qquad
\widetilde W_{(j)}^{b+1}
=\max\{2n_j,\widetilde W_{(j)}^b+\bar X^b_{(j)}-n_j\}.
\]
Then \(\widetilde W^b_{(j)}\ge \bar W^b_{(j)}\) for all \(b\ge0\) and \(j\in[m]\).
\end{lemma}
\noindent The proof is in Appendix~\ref{appendix_lemma::proof of multiple-type coupled dominace}. Let \(\mu_j=\lambda_j\Delta^\pi\). The increment in the dominating process is \(\bar X^b_{(j)}-n_j\), whose mean is \(\mu_j-n_j\le0\) by~\eqref{eq:policy_constraints}. If \(\mu_j=n_j\), the same random-walk maximum argument used in the single-type proof gives
\[
    \mathbb{E}[\bar W^B_{(j)}]
    \le \mathbb{E}[\widetilde W^B_{(j)}]
    =O(B^{1/2}).
\]
The same bound holds for every prefix \(b\le B\), so
\[
    \frac{1}{B}\sum_{b=0}^{B-1}\mathbb{E}[\bar W^b_{(j)}]=O(B^{1/2}).
\]
If \(\mu_j<n_j\), the dominating process has strictly negative drift. Since
\[
\left|\mathbb{E}[\bar X^b_{(j)}-n_j]\right|=n_j-\mu_j,
\]
by Kingman's bound \citep{kingman1962queues}:
\[
    \sup_{B\ge0}\mathbb{E}[\bar W^B_{(j)}]
    \le
    2n_j+
    \frac{\mathbb{E}[(\bar X^b_{(j)}-n_j)^2]}{2(n_j-\mu_j)}
    <\infty .
\]
Combining the zero-drift and negative-drift cases, under the nonpositive drift condition in~\eqref{eq:policy_constraints},
\[
    \sum_{j=1}^m \mathbb{E}[\bar W^B_{(j)}]=O(B^{1/2}),
    \qquad
    \frac{1}{B}\sum_{b=0}^{B-1}\sum_{j=1}^m \mathbb{E}[\bar W^b_{(j)}]=O(B^{1/2}).
\]
Under strict slack \(\mu_j<n_j\) for all \(j\), both quantities are \(O(1)\). If equality holds for at least one type and strict inequality holds for the others, the zero-drift types dominate the aggregate bounds, so the first part of the theorem applies. Types with zero arrival rates can be omitted from \([m]\).

It remains to translate these \(B\)-interval bounds into the theorem's finite-horizon metrics. The number of completed embedded full-threshold intervals in the physical horizon \([0,T]\) is exactly
\[
    B_\zeta(T)=\max\{b:u_b\le T\}
    =\left\lfloor\frac{\zeta T}{\Delta^\pi}\right\rfloor,
\]
and the residual time is less than one embedded full-threshold interval. This is the number of deterministic review intervals in the embedded process, not the number of actual event-driven WAIT batches. Since \(B_\zeta=\zeta T/\Delta^\pi+O(1)\), \(B_\zeta^{-1/2}=O((\zeta T)^{-1/2})\) and \(B_\zeta^{-1}=O((\zeta T)^{-1})\).

For throughput, terminal entry backlog controls the completion shortfall. More precisely, for the embedded process,
\[
    \throughput^*-\mathbb{E}[\throughput^{(\zeta,\bar\pi)}]
    \le
    \frac{1}{\zeta T}\sum_{j=1}^m l_j^\prime\,
    \mathbb{E}\!\left[
    \bar W^{B_\zeta}_{(j)}
    +A_j(u_{B_\zeta},T]
    +(l_j^\prime+2)n_j
    \right].
\]
The second term accounts for final partial-interval arrivals, and the last term absorbs the bounded in-service pipeline, startup effects, and the one-interval lag from Lemma~\ref{lem:wait_one_slot_dominance}. Since \(T-u_{B_\zeta}<\Delta^\pi/\zeta\), \(\mathbb{E}[A_j(u_{B_\zeta},T]]\le \lambda_j\Delta^\pi=O(1)\). Thus the terminal-backlog term gives an \(O((\zeta T)^{-1/2})\) gap under nonpositive drift and an \(O((\zeta T)^{-1})\) gap under strict slack.

For delay, terminal backlog is not the right object; the relevant quantity is the time average of the entry queues. Let \(\bar Q^b=\sum_{j=1}^m \bar W^b_{(j)}\). In physical time, each embedded full-threshold interval has length \(\Delta^\pi/\zeta\). The theorem, however, reports service-normalized delay, multiplying elapsed physical time by \(\zeta\). Thus each embedded interval contributes a constant scaled time \(\Delta^\pi\). The entry-wait integral in service-normalized units is bounded by
\[
    \Delta^\pi\sum_{b=0}^{B_\zeta-1}\bar Q^b+O(B_\zeta),
\]
where the \(O(B_\zeta)\) term accounts for the embedded-interval convention and bounded decode-stage waiting. The expected number of completed prompts is \(\Theta(B_\zeta)\): cumulative threshold-batch removals equal cumulative arrivals minus terminal backlog, and the missed removals are \(O(B_\zeta^{1/2})\) at the boundary and \(O(1)\) under strict slack. Therefore
\[
    \mathbb{E}[\Latency^{(\zeta,\bar\pi)}],
    \mathbb{E}[\TTFT^{(\zeta,\bar\pi)}]
    \le
    O\!\left(
    \frac{1}{B_\zeta}\sum_{b=0}^{B_\zeta-1}
    \mathbb{E}[\bar Q^b]
    \right)+O(1).
\]
The prefix bound gives
\[
    \mathbb{E}[\Latency^{(\zeta,\bar\pi)}],
    \mathbb{E}[\TTFT^{(\zeta,\bar\pi)}]
    =
    O(B_\zeta^{1/2})
    =
    O((\zeta T)^{1/2})
\]
under nonpositive drift. The corresponding calendar-time delay is smaller by a factor \(1/\zeta\). Under strict slack, the same service-normalized expression is \(O(1)\). Lemma~\ref{lem:wait_one_slot_dominance} transfers these throughput, latency, and TTFT bounds from \(\bar\pi\) to the actual event-driven WAIT policy, completing the proof of Theorem~\ref{thm:wait_heavy_traffic}.
